\documentclass[10pt,graphics,aspectratio=169,table]{beamer}
\usepackage{listings}
\usepackage{amsmath}
\usepackage{hyperref}

\usetheme{metropolis}

\title{Introduction to \TeX}
\author{Jakob Krebs}
\date{ESE 2019}
\institute{NERD101 - ESE - ifsr - TU Dresden}
\titlegraphic{\hfill\includegraphics[height=1.25cm]{../logo}}
\begin{document}
\maketitle
\begin{frame}{Outline}
    \tableofcontents
\end{frame}
\section{What ist \TeX?}
\begin{frame}{What is \TeX}
    \begin{itemize}
        \item Typesetting tool developed by Donald Knuth
        \item enables people to write Text and let the formatting be done by the computer
        \item it is the defacto standart for scientific writing
    \end{itemize}
\end{frame}
\section{Math mode}
\begin{frame}[fragile]{Math mode}
        \begin{equation}
            \hat{y} = \begin{pmatrix}y_1 \\ \vdots \\ y_n \end{pmatrix}, \hat{y'} = \begin{pmatrix}y'_1 \\ \vdots \\ y'_n \end{pmatrix}, A \in \mathbb{R}^{n \times n}
        \end{equation}
    Try, doing this in word, pages or libreoffice
    \begin{lstlisting}[language=tex]
\hat{y} = \begin{pmatrix}y_1 \\ \vdots \\ y_n \\
\end{pmatrix}, \hat{y'} = \begin{pmatrix}y'_1 \\ \vdots \\
y'_n \end{pmatrix}, A \in \mathbb{R}\^{n \times n}
\end{lstlisting}
\end{frame}

\section{a simple document}
\begin{frame}[fragile]{a simple document}
\begin{lstlisting}[language=tex]
\documentclass[12pt, letterpaper, twoside]{article}
\usepackage[utf8]{inputenc}
 
\title{this document}
\author{this author}
\date{ESE 2019}
 
\begin{document}
Lorem ipsum dolor sit. 
\end{document}
\end{lstlisting}
\end{frame}

\section{structuring a text}
\begin{frame}[fragile]{structureing a text}
    we can structure our text with sections and subsections.
    a table of contents will be created for us
    \begin{lstlisting}[language=tex]
    \chapter
    \section
    \subsection
    \subsubsection
    \tableofcontents
    \end{lstlisting}
\end{frame}

\section{presentations}
\begin{frame}[fragile]{presentations with beamer}
bla fasel laber
\end{frame}

\section{macros}
\begin{frame}[fragile]{macros}
No one whats to type this to create a matrix
    \begin{lstlisting}[language=tex]
    \begin{pmatrix} a & b \\ c & d \end{pmatrix}
    \end{lstlisting}
so we define a macro
    \begin{lstlisting}[language=tex]
\newcommand{\m}[1] {
    \begin{pmatrix}
    #1
    \end{pmatrix}}
    \end{lstlisting}
and now we can type
    \begin{lstlisting}[language=tex]
    \m{a&b\\c&d}
    \end{lstlisting}
\end{frame}

\section{reference management}
\begin{frame}[fragile]{reference management with bibtex}
    \begin{itemize}
        \item it can be hard to keep track of all sources one used
        \item referncing papers, articles or books can be hard
        \item different standarts on referencing are confusing
    \end{itemize}
    \begin{lstlisting}[language=tex]
@inbook{knuth,
    author = "Donald E. Knuth",
    title = "Fundamental Algorithms",
    publisher = "Addison-Wesley",
    year = "1973",
    chapter = "1.2",
    keywords  = "knuth,programming"
}
    \end{lstlisting}
\end{frame}

\begin{frame}[fragile]{including zour bibtex file}
    \begin{lstlisting}[language=tex]
\usepackage{biblatex}
\addbibresource{sample.bib}
    \end{lstlisting}

now you just reference it in your text with

    \begin{lstlisting}[language=tex]
\cite{knuth}
\end{lstlisting}
and print your references at the end of your document
\begin{lstlisting}
\printbibliography
\end{lstlisting}
\end{frame}

\section{useful tools and references}
\begin{frame}{useful things}
you will encounter problems like everywhere else
\begin{itemize}
    \item not knowing the command for a symbol \url{http://detexify.kirelabs.org}
    \item simple manuals on \TeX topics \url{https://www.overleaf.com/learn/latex}
    \item other problems $\implies$ the search engine of least distrust
\end{itemize}
\end{frame}
\end{document}
